\secciongrande{Introducción al desarrollo web en entorno cliente}

\resumen{
  La web surgió inicialmente como un medio de documentos hipertextuales basados en HTML, alojados en servidores y accedidos mediante navegadores. El consorcio W3C estandariza las tecnologías web. El desarrollo web es multidisciplinar, abarcando desde la captación de necesidades del cliente hasta la presentación final del sitio. El modelo cliente-servidor es la arquitectura dominante, donde el servidor almacena y el navegador presenta la información.
}

\definicion{WEB}{Sistema de información basado en hipertexto que permite la creación, publicación y consulta de documentos electrónicos a través de internet. La web se caracteriza por su capacidad de enlazar recursos mediante referencias (hipervínculos).}

\definicion{Hipertexto}{Modelo de presentación de información donde los elementos están interconectados mediante referencias que permiten acceder a información relacionada de forma no lineal. En la web, los hipervínculos permiten navegar entre documentos.}

\definicion{W3C}{World Wide Web Consortium (Consorcio Mundial de la Web). Organización encargada de desarrollar y mantener los estándares web, incluyendo HTML, CSS, XML, SVG y otros lenguajes y tecnologías que forman la base de internet. Fundado por Tim Berners-Lee en 1994.}

\subtituloA{Origen y evolución de la web}

\lista{
  **Modelo de servidor:** Las páginas se alojaban en servidores conectados 24 horas al día, 365 días al año, transmitiéndose a los equipos que solicitaban el contenido mediante protocolos como HTTP. |
  **Estandarización:** La explosión de la web llevó a la creación del W3C, que estableció estándares para garantizar la interoperabilidad entre diferentes navegadores y dispositivos.
}
